\documentclass[letter, 10pt]{article}
\usepackage[utf8]{inputenc}
\usepackage[spanish]{babel}
\usepackage{amsfonts}
\usepackage{amsmath}
\usepackage[dvips]{graphicx}
\usepackage{url}
\usepackage[top=3cm,bottom=3cm,left=3.5cm,right=3.5cm,footskip=1.5cm,headheight=1.5cm,headsep=.5cm,textheight=3cm]{geometry}



\begin{document}
\title{Inteligencia Artificial \\ \begin{Large}Estado del Arte: Problema [Nombre Problema]\end{Large}}
\author{[Nombre autor]}
\date{\today}
\maketitle


%--------------------No borrar esta secci\'on--------------------------------%
\section*{Evaluaci\'on}

\begin{tabular}{ll}
Resumen (5\%): & \underline{\hspace{2cm}} \\
Introducci\'on (5\%):  & \underline{\hspace{2cm}} \\
Definici\'on del Problema (10\%):  & \underline{\hspace{2cm}} \\
Estado del Arte (40\%):  & \underline{\hspace{2cm}} \\
Modelo Matem\'atico (20\%): &  \underline{\hspace{2cm}}\\
Conclusiones (15\%): &  \underline{\hspace{2cm}}\\
Bibliograf\'ia (5\%): & \underline{\hspace{2cm}}\\
 &  \\
\textbf{Nota Final (100\%)}:   & \underline{\hspace{2cm}}
\end{tabular}
%---------------------------------------------------------------------------%
\vspace{2cm}


\begin{abstract}
El Truck and Trailer Routing Problem (TTRP) consiste en la planificacion de rutas de reparto a multiples clientes con demandas distintas, en donde
se tiene una cantidad definida de camiones y cargamentos. Ciertos clientes pueden ser servidos por camiones con cargamento y otros exigen ser servidos
por camiones desacoplados (ya sea por razones geograficas o logisticas).

\end{abstract}

\section{Introducci\'on}
El presente informe tiene el objetivo de explicar el Truck and Trailer Routing Problem (TTRP), empezando por su definicion original y
posteriormente desglosando articulos clave para caracterizar el estado del arte del problema. Para cada articulo se hara un resumen breve,
la razon de su eleccion y su impacto en la forma de encontrar soluciones optimas.

Se definira el problema en relacion a sus variables, restricciones, funcion objetivo y parametros. Usando el articulo inicial como
base para luego detallar los variantes presentados en articulos posteriores, o de ser el caso, explicar como estos difieren a la hora
de abordar la solucion al problema. Al tener muchas aplicaciones reales, existen multiples versiones que agregan variables para
acercarlo mas a la realidad, el informe no pretende desviar mucho del problema original si no que enfocarse en la evolucion de las
heuristicas empleadas.

El estudio de este problema ha sido fundamental en la logistica de repartos masivos, ya que en muchos casos es necesario usar multiples
tipos de vehiculo para satisfacer la demanda. El problema original surge en parte por problemas encontrados en la industria lactea, en
donde los caminos rurales podian ser demasiado estrechos, haciendo dificil el transito de camiones con carga. Tambien se ha utilizado en el
sector urbano, particularmente en partes de Brasil donde la distribucion de bebidas y tabaco ocurrian en zonas poblacionalmente densas.

\section{Definici\'on del Problema}
El primer articulo en definir este problema fue "A tabu search method for the truck and trailer routing problem - I-Ming Chao" (Chao 2002),
el cual consistia en encontrar la ruta mas corta para una cantidad determinada de clientes, camiones y cargamentos.

Se conoce la cantidad de clientes desde el inicio, en donde cada uno tiene una demanda explicita desde el inicio, la cual puede ser distinta
entre si. Los vehiculos empleados para el reparto pueden tomar dos formatos: con carga (vehiculos completos) y sin carga, en donde los clientes
pueden ser satisfechos con alguno de esos tipos. La capacidad de los vehiculos y sus cargas es igual para todos, en donde se conoce la cantidad
dispuesta desde el inicio.

Lo que lo diferencia de un "vehicle routing problem" (VRP) tradicional es que aqui los vehiculos pueden desacoplarse para modificar el
abanico de clientes a los cuales este puede atender, estacionandose en la posicion de un cliente y posteriormente recorriendo una "sub-ruta",
para la cual el vehiculo esta obligado a volver antes de retomar su curso normal.

Se exige que la demanda de todos los clientes sea satisfecha, la cual debe ser atendida por un vehiculo una sola vez, osea que no se puede
repartir la demanda de un cliente entre multiples visitas. Los vehiculos tienen una capacidad definida, por ende se deben atender los clientes
dentro del presupuesto de capacidad inicial, para la demanda de cada uno de estos. Finalmente, si un vehiculo completo se desacopla, solo este
puede volver a recoger su carga inicial.

El objetivo es minimizar la distancia recorrida por la flota completa, la cual considera un punto de partida para todos los vehiculos a la cual
todos deben volver una vez que la demanda total sea satisfecha.

Al ser un problema que se construye a partir de un VRP (el cual es NP-hard) un TTRP tiene por lo menos la misma complejidad, por lo que un
algoritmo exacto que pueda resolver el problema de forma optima es poco probable. El primer articulo propone buscar una solucion factible
primero, definiendo un recorrido completo para cada vehiculo, compuesto por rutas con vehiculo con carga y sub-rutas con el vehiculo
desacoplado.

Las soluciones iniciales se construyen a partir de una "semilla", la cual define a que cliente atiende cada vehiculo primero,
los clientes posteriores se eligen iterativamente en base a los mas lejanos a las semillas anteriores y al deposito (se eligen los mas lejanos
primero). Sin embargo, antes de llegar a las soluciones factibles se implementa una relajacion al problema; se permite que los clientes sean
satisfechos de forma fraccionaria, para la cual se aplica un redondeo una vez que se asignen todos los clientes, reduciendo el costo computacional
del algoritmo. Una vez que se asignan los clientes a las rutas, se crea la ruta para clientes que pueden ser atendidos por vehiculos completos,
posteriormente, se analizan los clientes con rutas desacopladas, para los cuales se crea una sub-ruta para el mas cercano a la ruta principal,
si se pueden atender mas clientes, estos se agregan a la sub-ruta, si la capacidad no da o esta muy lejos, se crea otra sub-ruta.

Luego de construir las soluciones iniciales, se aplica el procedimiento "descent improvement", necesario para transformar nuestras rutas
iniciales en soluciones factibles (efecto de la relajacion aplicada).
El primer paso consiste en mover clientes entre las rutas propuestas, para el cual se define una penalizacion que se calcula a partir de
la demanda, la capacidad y el tipo del vehiculo. Se examinara un solo cliente a la vez, el cual no puede moverse si este se encuentra como
punto de estacionamiento o si la ruta destino no es compatible.
El segundo paso consiste en cambiar a dos clientes de dos rutas distintas, usando un sistema de penalizacion similar al punto anterior, en
donde los cambios mantienen la factibilidad en relacion a la capacidad del camion.
En el tercer paso se analizan las sub-rutas, para las cuales se eliminan las aristas que conectan al inicio de la sub-ruta (osea el inicio y
el final) y se remplazan por aquellas que conectan dos clientes consecutivos de la sub-ruta a un cliente en la ruta principal, en terminos
sencillos, se mantienen los clientes de la sub-ruta pero se conectan de forma distinta a la principal.

Para cada solucion inicial, el procedimiento anterior correra multiples veces hasta que ninguno de los pasos se pueda ejecutar, cuando esto
ocurra, se aplicara una limpieza de tipo "2-opt", una optimizacion conocida para problemas de grafos, la cual consiste en eliminar las rutas
"cruazadas" entre si (comunmente usando en problemas de tipo Traveling Salesperson Problem).

Los pasos anteriores generaran una serie de soluciones factibles pero no necesariamente optimas y es aqui donde entra la parte mas interesante
del articulo, se emplea una heuristica de mejora de tipo "tabu search", para la cual se desarrolla una variante nueva llamada "ojective-based tabu
restriction (OTB)". En terminos muy simples, tabu search consiste en mejorar una solucion inicial al desplazarse entre soluciones vecinas, para la
cual elige el vecino que de el valor mas optimo (incluso si es peor al original), para evitar ciclos esta define una tabla "tabu", que guarda los
caminos recientemente recorridos, permitiendo que este diversifique. Por otro lado, la heuristica nueva define un rango de mejora, en donde, cualquier
valor que se desvie de este rango estara prohibido, se evaluaran los vecinos dentro de un numero de iteraciones definidas y si no encuentra soluciones,
incrementara la desviacion permitida. La combinacion de ambas heuristicas hace que las iteraciones no repitan pasos ciclicamente y que ademas escape
de los maximos locales, ya que los caminos van mejorando gradualmente y se van prohibiendo a la vez, abriendo el paso para alcanzar soluciones
aun mejores. Finalmente, se aplica un criterio de aspiracion (AC) que excluya soluciones en la tabla tabu que son peores a la actual.

\section{Estado del Arte}
Debido al gran impacto que tiene este problema, han surgido multiples articulos abordando TTRP y sus variantes. Uno de los mas interesantes se encuentra
en "A note on the truck and trailer routing problem" (Lin, Yu. Chou 2010), el cual plantea una leve variante denominada "relaxed truck and trailer routing
problem" (RTTRP), el cual evalua la posibilidad de variar la cantidad de vehiculos inciales para asi reducir la distancia recorrida por la flota. El paper
hace un benchmark comparandolo directamente con la propuesta del 2002 y la mejor solucion de ese entonces obtenida el 2009, obteniendo resultados consistentemente
mejores que ambos, con un margen promedio de 5 porciento mejor que la ultima.

Años mas tarde, aparece "A matheuristic for the truck and trailer routing problem" (Villegas, Prins, Prodhon, Medaglia, Velasco 2013) con el fin de emplear
los ultimos metodos de optimizacion hibrida. Las matheuristicas surgen como la alternativa a los metodos exactos, las heuristicas y las metaheuristicas. Este
acercamiento descompone el problema inicial en multiples sub-espacios, restringiendo el espacio de busqueda a vecinos locales y usando programacion matematica
para hallar una mejor solucion. Se emplea un metodo hybrid GRASP x ILS como columnas en una formulacion de set-partitioning. Este nuevo metodo logro los mejores
resultados en terminos de calidad y tiempo de computo para su epoca, abordando tambien RTTP para flotas con tamano variable.

El año 2025 aparecio el articulo "A Cooperative Algorithm to Tackle the Truck and Trailer Routing Problem" el cual planteaba resolver el problema usando
algoritmos cooperativos, metodo que consiste en separar un algoritmo en multiples procesos independientes que trabajan en subconjuntos de informacion que
los otros procesos desconocen. Dichos espacios se resuelven por si solos y luego van "evolucionando", generando nuevos espacios de soluciones que sirven
para iterar y llegar a nuevos optimos. En particular, este paper emple el algorimo "CoopA", el cual es competitivo con otros del mismo tipo, por lo que
abre las puertas a una nueva area a estudiar.

Finalmente, uno de los ultimos articulos publicados del tema es "Nuevos optimos para el Problema de Rutas de Camiones y Remolques con Restricciones de Capacidad
usando un algoritmo colaborativo" (Navarro 2026), el cual combina multiples aspectos de los descubrimientos mencionados. Esta propuesta hace uso de CORAL, un algoritmo de busqueda
local centrado en resolver el CTTRP. CORAL consta de dos fases, la primera de contruccion, basado en procedimientos GRAASP y una fase posterior de reparacion basada
en procedimientos Hill-climbing. Esta propuesta obtiene optimos iguales a los de la literatura de forma consistente, pero al mismo tiempo siendo considerablemente
mas eficiente.

Es importante notar que hasta la fecha, el problema siempre se aborda con una metodologia de dos pasos: encontrar una solucion factible rapido y posteriormente
mejorarla con heuristicas especificas, evidente en todos los articulos discutidos hasta ahora.

\section{Modelo Matem\'atico}
En esta secci\'on se describe el modelo matem\'atico del problema planteado en el art\'iculo original del 2002, desglosando las variables, restricciones, funci\'on objetivo
y par\'ametros en cada secci\'on.

\subsection{Esquema General}

El siguiente esquema resume, de forma gr\'afica, los tres tipos de ruta que conforman una soluci\'on del TTRP: la ruta pura de cami\'on (s\'olo clientes t.c.), la ruta completa de cami\'on con remolque (tour principal sobre clientes v.c.) y el sub-tour que se genera al desacoplar el remolque en un punto de estacionamiento.

\begin{figure}[h]
\centering
\setlength{\unitlength}{1mm}
\begin{picture}(150,82)(0,0)
 \put(55,70){\framebox{\parbox{24mm}{\centering\footnotesize Dep\'osito}}}
 \put(67,60){\line(0,1){10}}
 \put(20,60){\line(1,0){94}}
 \put(20,48){\line(0,1){12}}
 \put(114,48){\line(0,1){12}}

 \put(0,30){\framebox{\parbox{38mm}{\centering\footnotesize Ruta pura de cami\'on\\ (s\'olo clientes t.c.)}}}
 \put(84,30){\framebox{\parbox{60mm}{\centering\footnotesize Ruta completa (cami\'on + remolque)\\ tour principal: clientes v.c.}}}

 \multiput(114,13)(0,2){7}{\line(0,1){1}}

 \put(84,0){\framebox{\parbox{60mm}{\centering\footnotesize Sub-tour: clientes t.c.\\ (remolque estacionado en el punto de desacople)}}}
\end{picture}
\caption{Esquema de los tres tipos de ruta del TTRP. Todas las rutas comienzan y terminan en el dep\'osito. L\'inea s\'olida: asignaci\'on entre el dep\'osito y un tipo de ruta. L\'inea punteada: desacople hacia un sub-tour, recorrido \'unicamente por el cami\'on mientras el remolque permanece estacionado.}
\end{figure}

\subsection{Par\'ametros}
Los par\'ametros del sistema se conocen desde un principio. La siguiente tabla resume la notaci\'on utilizada.

\begin{table}[h]
\centering
\begin{tabular}{cl}
\hline
Par\'ametro & Descripci\'on \\
\hline
$q_i$ & Demanda del cliente $i$ \\
$t_i$ & Tipo del cliente $i$ \\
$c_{ij}$ & Distancia eucl\'idea entre el cliente $i$ y el cliente $j$ \\
$Q_1$ & Capacidad de cada cami\'on \\
$Q_2$ & Capacidad de cada remolque \\
$m_1$ & N\'umero de camiones disponibles \\
$m_2$ & N\'umero de remolques disponibles ($m_1 > m_2$) \\
\hline
\end{tabular}
\caption{Par\'ametros del modelo.}
\end{table}

Primero se define la demanda de cada cliente, $q_i$: la demanda del cliente $i$.

En donde cada cliente tiene un tipo $t_i$, que restringe el tipo de veh\'iculo que puede atenderlo:
\begin{equation}
t_i =
\begin{cases}
0 & \text{vehicle customer (v.c.)} \\
1 & \text{truck customer (t.c.)}
\end{cases}
\end{equation}

Un v.c. puede ser atendido por cami\'on solo o por cami\'on + remolque.
Un t.c. solamente puede ser atendido por cami\'on solo.

Luego, tenemos la distancia desde un cliente ($i$) a otro ($j$):
\begin{equation}
c_{ij} = \text{distancia eucl\'idea entre } i \text{ y } j
\end{equation}

Para la flota, cada cami\'on y remolque tienen capacidades aparte e iguales para cada tipo: $Q_1$, capacidad de cada cami\'on; $Q_2$, capacidad de cada remolque.

Por lo que la capacidad de un veh\'iculo completo es $Q_1 + Q_2$.

La capacidad de la flota tambi\'en se conoce de antemano y es fija; para problemas de RTTRP esto pasa
a ser una variable, pero para el art\'iculo original es una constante: $m_1$, n\'umero de camiones; $m_2$, n\'umero de remolques.

En donde se establece que $m_1 > m_2$. La posici\'on del dep\'osito se omite aqu\'i por simplicidad.

\subsection{Variables}

De forma preliminar, se separan los clientes por ruta mediante la variable de asignaci\'on $x_{ij}$:
\begin{equation}
x_{ij} =
\begin{cases}
1 & \text{si el cliente } i \text{ pertenece a la ruta } j \\
0 & \text{si no}
\end{cases}
\qquad i=1,\dots,n, \quad j=1,\dots,m_1
\end{equation}

El art\'iculo utiliza expl\'icitamente variables $x_{ij}$ para su problema de asignaci\'on relajado.

Luego, para cada ruta hay que definir lo siguiente:

\begin{enumerate}
\item Cu\'ales son los puntos de acople y desacople.
\item Qu\'e clientes van en un sub-tour.
\item En qu\'e orden se visitan.
\item Cu\'antos sub-tours hay.
\end{enumerate}

\subsection{Restricciones}

\begin{enumerate}
\item Todos los clientes deben ser atendidos:

Cada cliente $i$ debe pertenecer a una ruta.

\begin{equation}
\sum_{j=1}^{m_1} x_{ij} = 1, \qquad i=1,\dots,n
\end{equation}

El art\'iculo utiliza precisamente esta restricci\'on en su problema de asignaci\'on.

\item No se puede dividir la demanda.

Esta restricci\'on va de la mano con la anterior ya que si pertenece a una ruta, se fuerza a que se satisfaga una sola vez.

\item Se respetan las capacidades.

La demanda total de una ruta completa (cami\'on + remolque, incluyendo el tour principal y sus sub-tours) debe cumplir:

\begin{equation}
\sum_{i=1}^{n} q_i \, x_{ij} \le Q_1 + Q_2, \qquad j=1,\dots,m_2
\end{equation}

Una pure truck route debe cumplir:

\begin{equation}
\sum_{i=1}^{n} q_i \, x_{ij} \le Q_1, \qquad j=m_2+1,\dots,m_1
\end{equation}

Lo mismo ocurre con un sub-tour:

\begin{equation}
\sum_{i \,\in\, \text{sub-tour}} q_i \le Q_1
\end{equation}

(con $i$ siendo del sub-tour) porque durante el sub-tour est\'a circulando solamente el cami\'on.

\item Tipos de cliente.

Se debe respetar que cada tipo de cliente sea atendido por el veh\'iculo del tipo correspondiente.
\end{enumerate}

\subsection{Funci\'on Objetivo}

Aqu\'i el objetivo se define simplemente como minimizar la distancia total recorrida por la flota. Para las conexiones utilizadas $(i,j)$, cada una tiene distancia $c_{ij}$, conceptualmente:

\begin{equation}
\min \sum_{i=1}^{n} \sum_{j=1}^{m_1} c_{ij} \, x_{ij}
\end{equation}

La formulaci\'on concreta del art\'iculo utiliza la distancia total recorrida como objetivo, asumiendo que el costo es proporcional a la distancia.

\section{Conclusiones}
Se puede concluir que TTRP es un problema de suma importancia, el cual tiene activo desarrollo en la comunidad academica. El ultimo articulo empleado en este
informe es de hace solo un par de meses atras (desde la redaccion de este documento), por lo que se espera ver mas avances futuros. El objetivo de este estudio
era mostrar la evolucion de la solucion del problema, enfocandose en las estrategias y heuristicas empleadas, evitando tambien cubrir modificaciones demasiado
grandes al problema original. Las diferencias principales de cada hallazgo recaen en el metodo de optimizacion luego de la fase de construccion inicial, las
cuales hacen uso de los ultimos descubrimientos de la epoca (como matheuristicas u algoritmos cooperativos). Se espera que los siguientes avances sigan haciendo
uso de algoritmos cooperativos, ya que la computacion paralela es uno de los modelos mas en voga actualmente. Una propuesta interesante seria implementar mas tipos
de clientes para cada tipo de desacople, ya que actualmente el problema opera de forma binaria. Otro factor de alto impacto seria integrar variables de velocidad
y tiempo, para acercar mas el problema a la vida real. De todas formas, se espera que el problema se siga estudiando por mucho tiempo mas.

\section{Uso LLMs}

La redacci\'on del presente informe (todo el texto explicativo y los an\'alisis) fue realizada \'unicamente por el autor. Las herramientas de IA descritas a continuaci\'on se utilizaron exclusivamente como apoyo t\'ecnico y de investigaci\'on, nunca para generar o reemplazar la redacci\'on del contenido.

\begin{enumerate}

\item Herramienta utilizada: Claude Code (Claude Sonnet 5), plan de pago.

\item Prop\'osito de uso: Dar formato LaTeX a la secci\'on de Modelo Matem\'atico (tabla de par\'ametros, ecuaciones formales y un esquema de los tipos de ruta del TTRP), corregir un error de compilaci\'on preexistente (codificaci\'on del archivo e inconsistencias en una f\'ormula), corregir ortograf\'ia/tildes, y generar las entradas BibTeX de \texttt{Referencias.bib} a partir de la lectura de los PDF de los art\'iculos citados.

\item Prompts relevantes: ``mejora la secci\'on de modelo matem\'atico usando esquemas, gr\'aficos y f\'ormulas con formato latex'' (con la restricci\'on expl\'icita de no modificar la redacci\'on existente ni otros archivos sin autorizaci\'on); ``lee los pdf enumerados ... y agr\'egalos al archivo de referencias''.

\item Nivel de edici\'on: No se copi\'o y peg\'o texto explicativo generado por la IA; su uso se restringi\'o a c\'odigo LaTeX (tablas, ecuaciones, diagramas) y a metadatos bibliogr\'aficos, los cuales fueron revisados por el autor antes de aceptarlos.

\item Validaci\'on: El modelo matem\'atico formalizado se contrast\'o directamente contra la formulaci\'on expl\'icita del art\'iculo original (Chao, 2002) antes de incorporarlo; las referencias bibliogr\'aficas se generaron leyendo directamente los PDF originales de cada art\'iculo.

\item Herramienta utilizada: ChatGPT, con \emph{Retrieval-Augmented Generation} (RAG).

\item Prop\'osito de uso: Resolver dudas t\'ecnicas puntuales sobre el TTRP y sus variantes, aliment\'andolo directamente con los papers citados e incluidos en este proyecto para que las respuestas estuvieran ancladas a su contenido.

\item Prompts relevantes: Preguntas t\'ecnicas dirigidas sobre heur\'isticas, notaci\'on y resultados espec\'ificos de los art\'iculos cargados como contexto (RAG), usadas \'unicamente para aclarar conceptos antes de redactar.

\item Nivel de edici\'on: No se copi\'o ni peg\'o texto generado; las respuestas se usaron solo como apoyo de comprensi\'on, y la redacci\'on final es propia.

\item Validaci\'on: Al estar las respuestas ancladas (RAG) a los propios papers citados, se verificaron contra esos mismos documentos originales.

\end{enumerate}


\section{Bibliograf\'ia}

Todos los art\'iculos y la tesis citados en este informe fueron descargados utilizando el acceso institucional de la Universidad T\'ecnica Federico Santa Mar\'ia: las referencias \cite{Chao2002}, \cite{LinYuChou2010} y \cite{Villegas2013} se obtuvieron desde ScienceDirect (Elsevier); la referencia \cite{PerezRodriguezUrban2025}, desde el sitio de la revista \emph{Computaci\'on y Sistemas}; y la referencia \cite{Navarro2026}, desde el repositorio de tesis de la UTFSM.

\begin{thebibliography}{9}

\bibitem{Chao2002}
I.-M. Chao, ``A tabu search method for the truck and trailer routing problem,'' \emph{Computers \& Operations Research}, vol.~29, no.~1, pp.~33--51, 2002. doi: 10.1016/S0305-0548(00)00056-3.

\bibitem{LinYuChou2010}
S.-W. Lin, V. F. Yu, and S.-Y. Chou, ``A note on the truck and trailer routing problem,'' \emph{Expert Systems with Applications}, vol.~37, no.~1, pp.~899--903, 2010. doi: 10.1016/j.eswa.2009.06.077.

\bibitem{Villegas2013}
J. G. Villegas, C. Prins, C. Prodhon, A. L. Medaglia, and N. Velasco, ``A matheuristic for the truck and trailer routing problem,'' \emph{European Journal of Operational Research}, vol.~230, no.~2, pp.~231--244, 2013. doi: 10.1016/j.ejor.2013.04.026.

\bibitem{PerezRodriguezUrban2025}
R. P\'erez-Rodr\'iguez and L. E. Urb\'an-Rivero, ``A cooperative algorithm to tackle the truck and trailer routing problem,'' \emph{Computaci\'on y Sistemas}, vol.~29, no.~1, pp.~365--375, 2025. doi: 10.13053/CyS-29-1-4842.

\bibitem{Navarro2026}
A. F. Navarro Galleguillos, ``Nuevos \'optimos para el Problema de Rutas de Camiones y Remolques con Restricciones de Capacidad usando un algoritmo colaborativo,'' Tesis de Mag\'ister en Ciencias de la Ingenier\'ia Inform\'atica, Universidad T\'ecnica Federico Santa Mar\'ia, Departamento de Inform\'atica, Valpara\'iso, Chile, 2026.

\end{thebibliography}

\end{document} 
